\slajdid{02-s046}
\begin{frame}[label=02-s046]
\gusttekst
Međutim sa stanovišta broja čipova, situacija može biti značajno drugačija ili ista ili \ldots\ Za realizaciju u obliku zbira proizvoda
\[
F=\bar D\bar B+D\bar C\bar A
\]
treba nam 4 invertora da bi napravili \(\bar D,\bar C,\bar B\ i\ \bar A\) što ima u 1 čipu (ima ih 6 u 14pinskom pakovanju, 14-2 za masu i napajanje = 12, dva po invertoru, 12/2=6) jedno dvoulazno I kolo što ima u 1 čipu, jedno troulazno I kolo takođe u 1 čipu i na kraju 1 čip sa dvoulaznim ILI kolima što je ukupno 4 čipa. Za realizaciju u obliku proizvoda zbirova
\[
F=(D+\bar B)(\bar D+\bar C)(\bar D+\bar A)
\]
treba nam 4 invertora da bi napravili \(\bar D,\bar C,\bar B\ i\ \bar A\) što ima u 1 čipu, tri dvoulazna ILI kolo što ima u 1 čipu (ima ih 4 u 14pinskom pakovanju, 14-2 za masu i napajanje = 12, tri po dvoulaznom kolu, 12/3=4) i jedno troulazno I kolo takođe u 1 čipu što je ukupno 3 čipa. Sa ovog stanovišta je povoljnija realizacija u obliku proizvoda zbirova.
\end{frame}
